\chapter{Introduction}

\section{Background}
Modern wireless communication systems operate in environments where the
transmitter, channel, and interference conditions may change continuously.
An intelligent receiver must therefore perform more than symbol detection
under a fixed configuration. It should be able to observe the received
waveform, infer the signal format, select an appropriate processing chain,
and report the reliability of the recovered information. Automatic modulation
classification (AMC) is a central function in this adaptive receiver
architecture.

In a conventional receiver, the modulation format is known in advance. The
receiver can consequently configure a matched demodulator, synchronizer, and
decoder before data reception begins. In a non-cooperative or dynamically
configured system, this assumption is not valid. The receiver may encounter
BPSK, QPSK, M-PSK, or QAM waveforms without receiving an explicit
configuration message. AMC addresses this problem by estimating the
modulation format from the received signal.

\section{Motivation}
Classical AMC algorithms use statistical moments, cyclostationary features,
likelihood tests, or constellation-derived rules. These methods are useful
and interpretable, but their performance can deteriorate when the signal is
weak, the channel is fading, or the receiver has frequency and phase offsets.
Deep learning can learn discriminative representations directly from I/Q
samples and can be integrated with a software-defined receiver
\cite{oshea2016convolutional}.

The Phase 1 work established a raw-I/Q CNN baseline using RadioML 2016.10a.
However, a classification score alone does not demonstrate that the receiver
can successfully recover transmitted information. Phase 2 therefore moves
from an offline classification problem to a communication-system problem.
The predicted modulation is used as a control signal for adaptive
demodulation, and the final evaluation is performed in terms of BER, SER,
packet success, and throughput.

\section{Problem Statement}
The problem addressed in this report is the design and evaluation of an
intelligent adaptive digital communication receiver that can:
\begin{enumerate}
\item identify the modulation format of a received signal;
\item operate under noise, fading, phase offset, and frequency offset;
\item select a compatible demodulation method automatically; and
\item quantify the effect of modulation uncertainty on recovered information.
\end{enumerate}

\section{Aim of Phase 2}
The aim is to develop a reproducible transmitter--channel--receiver
simulation that demonstrates how AMC can be used as a communication control
function rather than as an isolated classifier. The project is intentionally
software-based so that every channel parameter and experiment can be repeated
before future validation with an SDR platform.

\section{Objectives}
\begin{enumerate}
\item Retain the Phase 1 RadioML CNN as a measured baseline.
\item Implement BPSK, QPSK, 8-PSK, and 16-QAM transmitters and demodulators.
\item Model AWGN, block-Rayleigh fading, carrier phase offset, and frequency
      offset.
\item Implement synchronization, channel equalization, modulation selection,
      and adaptive demodulation.
\item Measure BER, SER, packet success, throughput, and decision accuracy over
      a range of SNR values.
\item Compare the adaptive receiver with a fixed-QPSK receiver.
\item Prepare a reproducible and professionally documented Phase 2 report.
\end{enumerate}

\section{Contributions}
The main contribution is a communication-domain extension of the Phase 1
classifier. The work connects the signal-recognition stage to the physical
layer data path. It also introduces a multi-metric evaluation framework:
classification accuracy is complemented by BER, SER, packet success, and
effective throughput. This distinction is important because a correct
modulation decision is useful only when it leads to reliable bit recovery.

\section{Organization of the Report}
Chapter 2 reviews the communication principles and mathematical models used
in the project. Chapter 3 presents the proposed system architecture and
algorithms. Chapter 4 describes the software implementation, Phase 1
training, and Phase 2 experiment configuration. Chapter 5 presents the
measured results and comparison with Phase 1. The final chapter summarizes
the work, discusses limitations, and identifies future extensions.
